% from ega0-21-fr.tex
\oldpage[0\textsubscript{IV-1}]{170}
\begin{theorem}[21.7.1]
\label{0.21.7.1-fr}
\textnormal{(Cartier).}
Soient $K$ un corps, $L$ une extension de type fini de $K$. Alors
$\Omega^1_{L/K}$ et $\Upsilon_{L/K}$ sont des $L$-espaces vectoriels de rang
fini, et l'on a
\begin{equation}
\tag{21.7.1.1}
\label{0.21.7.1.1-fr}
\operatorname{rg}\Omega^1_{L/K}-\operatorname{rg}\Upsilon_{L/K}
=\deg.\operatorname{tr}_K L.
\end{equation}
\end{theorem}

\begin{proof}
Si $L'$ est un corps tel que $K\subset L'\subset L$, on a la suite exacte
(20.6.15.1) (appliquée à $\Lambda=P$, corps premier de $K$, $A=K$,
$B=L'$, $C=L$)
\[
0\to\Upsilon_{L'/K}\otimes_{L'}L\to\Upsilon_{L/K}
\to\Omega^1_{L'/K}\otimes_{L'}L\to\Omega^1_{L/K}
\to\Omega^1_{L/L'}\to0
\]
d'où (0$_{\rm III}$, 11.10.2)
\[
\operatorname{rg}_L\Omega^1_{L/K}-\operatorname{rg}_L\Upsilon_{L/K}
=(\operatorname{rg}_{L'}\Omega^1_{L/L'}-\operatorname{rg}_{L'}\Upsilon_{L/L'})
+(\operatorname{rg}_{L'}\Omega^1_{L'/K}-\operatorname{rg}_{L'}\Upsilon_{L'/K}).
\]
Comme par ailleurs $\deg.\operatorname{tr}_K L=\deg.\operatorname{tr}_K L'
+\deg.\operatorname{tr}_{L'}L$, on voit, par récurrence sur le nombre de
générateurs de l'extension $L$, que l'on est ramené à prouver (21.7.1.1)
lorsque $L=K(x)$. Distinguons alors trois cas~:
\begin{enumerate}
  \item[\rm a)] $x$ est transcendant sur $K$~; comme $L$ est séparable sur
  $K$, on a alors $\Upsilon_{L/K}=0$ (20.6.3)~; d'autre part, (20.5.9) et
  (20.4.13) montrent que $\Omega^1_{L/K}$ est de rang $1$, d'où
  (21.7.1.1) dans ce cas.
  \item[\rm b)] $L$ est une extension algébrique séparable de $K$, de sorte
  que l'on a encore $\Upsilon_{L/K}=0$ (20.6.3). D'autre part, on a
  $\Omega^1_{L/K}=0$ par (20.6.20), d'où encore (21.7.1.1).
  \item[\rm c)] $L$ est une extension algébrique inséparable de $K$~; le
  raisonnement du début montre que l'on peut se borner au cas où $L^p\subset K$~;
  il résulte donc de (21.4.8) que l'on a
  $\operatorname{rg}_L\Upsilon_{L/K}=\operatorname{rg}_L\Omega^1_{L/K}$,
  d'où encore (21.7.1.1).
\end{enumerate}
\end{proof}

\begin{corollary}[21.7.2]
\label{0.21.7.2-fr}
Soient $K$ un corps, $L$ une extension de type fini de $K$, $k$ un sous-corps
de $K$. Alors $\Omega^1_{L/K}$ et $\Upsilon_{L/K/k}$ sont des espaces
vectoriels de rang fini sur $L$, et l'on a
\begin{equation}
\tag{21.7.2.1}
\label{0.21.7.2.1-fr}
\operatorname{rg}\Omega^1_{L/K}-\operatorname{rg}\Upsilon_{L/K/k}
=\deg.\operatorname{tr}_K L+d(L/K,k).
\end{equation}
Par suite, on a l'inégalité
\begin{equation}
\tag{21.7.2.2}
\label{0.21.7.2.2-fr}
\operatorname{rg}\Omega^1_{L/K}-\operatorname{rg}\Upsilon_{L/K/k}
\geqslant\deg.\operatorname{tr}_K L.
\end{equation}
En outre, pour que les deux membres de (21.7.2.2) soient égaux, il faut et
il suffit que $k$ soit un corps admissible pour l'extension $L$ de $K$.
\end{corollary}

\begin{proof}
En effet, comme l'homomorphisme
$\Upsilon_{L/K}\to\Upsilon_{L/K/k}$ est surjectif, on a par définition
(21.6.1.2)
$\operatorname{rg}\Upsilon_{L/K}-\operatorname{rg}\Upsilon_{L/K/k}
=d(L/K,k)$, et le corollaire résulte donc aussitôt de (21.7.1) et de
(21.6.2).
\end{proof}

\begin{corollary}[21.7.3]
\label{0.21.7.3-fr}
Soient $k\subset K\subset L$ trois corps tels que $L$ soit une extension de
type fini de $k$. Alors on a
\begin{equation}
\tag{21.7.3.1}
\label{0.21.7.3.1-fr}
\operatorname{rg}_L\Omega^1_{L/k}-\operatorname{rg}_K\Omega^1_{K/k}
=\deg.\operatorname{tr}_K L+d(L/K,k).
\end{equation}
